\section{The remaining characteristic-three symplectic involutions}
\label{sec:symplectic-three-involutions}

\paragraph{Status.}
This separate complete derivation has passed whole-proof reviews by
/root and /root/imprimitive, with an independent exact supplement. It leaves D-HOS-13--16 unchanged. The new result
does not exclude all of the remaining natural ranks $4\le n\le283$.
It removes their projective involutions as a possible low-support
inertia class and specifies the resulting necessary class gate.
As before $\nu(T)$ is the codimension of the largest complex eigenspace.

\begin{lemma}[D-HOS-17: symmetric forms fixed by a complex structure]
\label{D-HOS-17}
Let $N=2m\ge4$, let $A\in\operatorname{GL}_{N}(3)$ satisfy
$A^2=-I$, and let $\mathcal O$ be any nonzero congruence orbit of
the general linear group on symmetric bilinear forms on
$\mathbf F_3^N$. Then
\[
 \frac{|\{B\in\mathcal O:A^tBA=B\}|}{|\mathcal O|}\le\frac1{39}.
 \tag{HOS3I.1}
\]
For rank two, the fraction is either zero or exactly
$2/[3(3^{N-1}-1)]$. For rank at least four it is at most $1/351$.
\end{lemma}

\begin{proof}
Put $\mathbf F_9=\mathbf F_3[t]$, $t^2=-1$, and let $t$ act as
$A$. If a symmetric form $B$ satisfies the fixed equation, then
$B(Ax,y)=-B(x,Ay)$, and
\[
                   H(x,y)=B(x,y)+tB(x,Ay)                    \tag{HOS3I.2}
\]
is Hermitian, linear in its first variable. Conversely the coefficient
of $1$ in a Hermitian form gives such a $B$. Their radicals agree;
an $\mathbf F_9$-rank-$r$ Hermitian form therefore gives an
$\mathbf F_3$-rank-$2r$ symmetric form. On the nondegenerate
quotient every Hermitian form is equivalent to $I_r$. Its underlying
symmetric form is $I_{2r}$: for $x=a+bt$ and $y=c+dt$, the
coefficient of $1$ in $x\bar y$ is $ac+bd$. Consequently only
the orthogonal type $\epsilon=(-1)^r$ occurs in rank $2r$.

We use the ordinary finite-form classification and group orders,
with their proofs in Taylor, \cite[pp.~115--118 and 138--141]{HOS:Taylor1992}.
These are orders of the full unitary and orthogonal groups, not their
special or projective quotients. The numbers of nondegenerate
Hermitian forms on $\mathbf F_9^r$ and symmetric forms of type
$\epsilon=(-1)^r$ on $\mathbf F_3^{2r}$ are respectively
\[
 \begin{split}
 H_r&=\frac{|\operatorname{GL}_r(9)|}{|\operatorname U_r(3)|}
      =3^{r(r-1)/2}\prod_{j=1}^{r}(3^j+(-1)^j),\\
 S_r&=\frac{|\operatorname{GL}_{2r}(3)|}
                  {|\operatorname O_{2r}^{(-1)^r}(3)|}
      =\frac12\,3^{r^2}(3^r+(-1)^r)
                       \prod_{j=1}^{r}(3^{2j-1}-1).
 \end{split}                                                   \tag{HOS3I.3}
\]
The formulas follow by orbit--stabilizer after choosing an ordered
basis; the orders in question are obtained in the cited source by
counting orthogonal or hyperbolic bases. Thus the relevant quotient is
\[
 R_r:=\frac{H_r}{S_r}
   =2\,3^{-r(r+1)/2}
       \frac{\prod_{j=1}^{r-1}(3^j+(-1)^j)}
            {\prod_{j=1}^{r}(3^{2j-1}-1)}.                     \tag{HOS3I.4}
\]
In particular $R_1=1/3$, $R_2=1/351$, and
\[
 \frac{R_{r+1}}{R_r}
     =3^{-r-1}\frac{3^r+(-1)^r}{3^{2r+1}-1}<1.               \tag{HOS3I.5}
\]

For a fixed rank, choose the radical first, then a nondegenerate form
on the quotient. An $A$-stable radical is an $\mathbf F_9$-subspace.
Hence the fraction in (HOS3I.1), when its type can occur, is exactly
\[
       \frac{{m\brack r}_{9}}{{2m\brack2r}_{3}}\,R_r,         \tag{HOS3I.6}
\]
where the brackets denote counts of subspaces. The first quotient is
at most one, since these $\mathbf F_9$-subspaces form a subset of
all the indicated $\mathbf F_3$-subspaces. For $r\ge2$ the bound
$R_r\le1/351$ proves the assertion. For $r=1$ the elementary
subspace counts give
\[
 \frac{{m\brack1}_{9}}{{2m\brack2}_{3}}\,R_1
     =\frac{2}{3(3^{N-1}-1)}\le\frac1{39},                   \tag{HOS3I.7}
\]
as $N\ge4$. Odd ranks and the other even-rank type have no fixed
forms. This proves the complete orbit statement.
\end{proof}

\begin{theorem}[D-HOS-18: the square-minus-one support bound]
\label{D-HOS-18}
Let $S=\operatorname{Sp}_{2n}(3)$, $n\ge4$, and let $V$ be any
nontrivial ordinary irreducible complex $S$-module of degree $a$.
If $h\in S$ satisfies $h^2=-I$, then
\[
 \nu(h|V)\ge\frac{19a}{78}\quad(n\text{ even}),\qquad
 \nu(h|V)>\frac{19a}{104}\quad(n\text{ odd}).                 \tag{HOS3I.8}
\]
The statement includes characters inflated from a represented
central quotient and makes no selfduality assumption.
\end{theorem}

\begin{proof}
Let $W$ be the natural space and $\langle\ ,\ \rangle$ its
alternating form. Regard $W$ as an $n$-dimensional space over
$\mathbf F_9$ by letting $t$ act as $h$. Then
\[
 \mathcal H(x,y)=\langle x,hy\rangle-t\langle x,y\rangle       \tag{HOS3I.9}
\]
is a nondegenerate Hermitian form, linear in the first variable.
For example $\mathcal H(hx,y)=t\mathcal H(x,y)$, by
$h^2=-I$ and symplecticity; the other identities follow directly.
The finite Hermitian Witt decomposition has $\lfloor n/2\rfloor$
hyperbolic planes and an anisotropic remainder of dimension zero
or one over $\mathbf F_9$.

Choose their two opposite maximal isotropic subspaces $L,L'$.
Their $\mathbf F_3$-dimension is
\[
              N=n\quad(n\text{ even}),\qquad
              N=n-1\quad(n\text{ odd}).                     \tag{HOS3I.10}
\]
They are $h$-stable, and $W_0=L\oplus L'$ is a nondegenerate
symplectic space of dimension $2N$ over $\mathbf F_3$.
Its orthogonal complement is zero or a symplectic plane. Let
$S_0=\operatorname{Sp}(W_0)$ act identically on that complement.
Inside $S_0$ take the opposite Siegel radicals $U,U^-$ associated
to $L,L'$. They are conjugate in $S_0$ and generate $S_0$.

Here generation can also be seen directly. The two radicals contain
the transvections with centers in $L$ and in $L'$. If
$z=x+y$, $x\in L$, $y\in L'$, and $\langle x,y\rangle\ne0$,
conjugation by a suitable transvection with center $x$ gives one
with center $z$. If this pairing vanishes and $y\ne0$, choose
$x'\in L$ with $\langle x',y\rangle\ne0$. First obtain the
center $x+x'+y$, then remove $x'$ by another such conjugation.
The cases $x=0$ or $y=0$ are already included. Thus all symplectic
transvections occur; elementary symplectic elimination shows that
they generate $S_0$, as also used in D-HOS-13.

Put $b=\dim V^{S_0}$ and $v_0=\dim V^U=\dim V^{U^-}$.
The intersection of these two invariant spaces is $V^{S_0}$,
so the subspace dimension formula gives
\[
                         v_0\le\frac{a+b}{2}.                \tag{HOS3I.11}
\]
When $n$ is even, $S_0=S$ and $b=0$. When $n$ is odd, $n\ge5$
and $2N=2n-2\ge8$. This is exactly a restriction subgroup in
Larsen--Shalev--Tiep, \cite[Proposition~4.2.3, pp.~1909--1912]{HOS:LST2011}:
its natural space is a nondegenerate summand and the group fixes
the orthogonal complement pointwise. Applied to our nontrivial
irreducible representation of the full simply connected group, its
explicit estimate gives
\[
 \frac ba<3^{8/3-(2N)/2}=3^{11/3-n}
             \le3^{-4/3}<\frac14.                            \tag{HOS3I.12}
\]
The last inequality is equivalent to $81>64$. The complete cited
proof, in particular its symplectic step 4 and summation step 6,
was read; no equality of ordinary and projective group orders is
being used. Thus $v_0\le a/2$ for even $n$ and $v_0<5a/8$ for
odd $n$.

Identify $U$ with the additive symmetric matrices on $L$. Under
the nondegenerate matrix-trace pairing over $\mathbf F_3$, its
complex character group is the same symmetric-matrix space.
The Levi subgroup $\operatorname{GL}(L)\subset S_0$ acts by
congruence. Therefore all weight spaces in one such orbit have
equal multiplicity in the restriction of $V$. The element $h$
normalizes $U$; its restriction $A$ to $L$ satisfies $A^2=-I$,
and its fixed weights are precisely the forms in D-HOS-17.
Replacing $A$ by its inverse in the contragredient convention
does not change the fixed equation.

Consequently at most $1/39$ of the total nonzero weight-space
dimension is fixed by $h$. On the other weights $h$ acts in pairs,
since its square acts trivially by conjugation on $U$. For each
pair of equal-dimensional weight spaces, an eigenvector of $h$
is determined by its component in one space. Every eigenvalue
therefore occupies at most half of the pair's total dimension.
This remains valid whether the central involution acts as $1$
or $-1$ on $V$. Summing over pairs proves
\[
 \nu(h|V)\ge\frac12\frac{38}{39}(a-v_0).                    \tag{HOS3I.13}
\]
Combining (HOS3I.11)--(HOS3I.12) gives exactly (HOS3I.8).
The argument uses the actual irreducible representation and its
weight spaces, not only a character-ratio bound.
\end{proof}

\begin{corollary}[D-HOS-19: a necessary residual inertia class]
\label{D-HOS-19}
Consider the main sixfold source problem with a finite inducing
image having an irreducibly acting, transitive normal central
product of factors whose simple quotients are
$\operatorname{PSp}_{2n}(3)$, $n\ge4$. Suppose its factor modules
are nonselfdual. If its actual six-dimensional multiplicity
variation has period rank nine, then $4\le n\le283$, and at
least one genuine inertia has all of the following properties:
it preserves every tensor factor; it is not a scalar on the full
space; and in every nonscalar factor component, every noncentral
power of prime order modulo the symplectic center has odd prime
order $\ell\ge5$ and, in its odd-order lift, has no natural
$1$ or $-1$ eigenvalue. In particular the projective orders of
those noncentral factor components are coprime to six.

The final coprimality assertion concerns projective factor orders.
It does not impose that condition on an independent scalar part
of the actual inertia. No existence claim is made for these classes.
\end{corollary}

\begin{proof}
The rank interval is D-HOS-16. For nonselfdual modules the full
compatible normalizer of one factor is scalar times inner:
the only nonsquare diagonal outer class is class-inverting by
Vinroot and hence cannot stabilize the module. This is the exact
normalizer result already checked in D-HOS-15, including lifts
through represented central quotients.

Taking a power and multiplying by a central sign cannot increase
projective support. The elementary cases in D-HOS-13--14 give
support at least $2a/9$ for actual order three or $h^2=I$, and
at least $4a/15$ for odd prime order with a natural fixed plane.
D-HOS-18 now supplies the remaining projective-involution case
$h^2=-I$. Thus a nonscalar factor component with at least one
prime-order power outside the residual class in the statement
has support at least
\[
 \beta_n a,\qquad
 \beta_n=\frac29\ (n\text{ even}),\qquad
 \beta_n=\frac{19}{104}\ (n\text{ odd}).                      \tag{HOS3I.14}
\]
Convolution of tensor-factor spectra gives the same proportional
bound on the full degree $e$ whenever one component has this
property. A nontrivial tensor-factor permutation has support
at least $e/4$, and a nonidentity scalar moves all $e$ dimensions.

If no inertia had all the asserted residual properties, every
genuine inertia would therefore move at least $\beta_n e$.
The elementary degree bound D-HOS-02 gives $e\ge3^{n-1}$.
For the actual intermediate-cover variation,
$6=2e(h-1)+D$ rules out $h\ge2$, and $h=1$ permits at most
six genuine marks and period rank at most six. Rank nine thus
requires $h=0$, at least twelve genuine marks, and $D=2e+6$.
It follows that
\[
 e\le\frac6{12\beta_n-2}
   =9\ (n\text{ even}),\qquad
 e\le\frac{156}{5}\ (n\text{ odd}).                          \tag{HOS3I.15}
\]
These contradict respectively $e\ge27$ for even $n\ge4$ and
$e\ge81$ for odd $n\ge5$. Hence the required inertia exists.
If a noncentral factor's projective order were divisible by two
or three, taking a suitable power would give one of the excluded
prime-order cases. This proves the projective coprimality claim.
The argument uses the actual inducing multiplicity variation;
its coefficient field need not equal the original imaginary
quadratic field.
\end{proof}

\paragraph{Exact supplement and limits.}
The accompanying exact finite reconstruction
enumerates every symmetric form in dimensions two and four,
classifies its rank and type by exact finite-field elimination,
and compares the fixed counts with (HOS3I.6). It also checks the
Hermitian construction, radical matrices and exact rational
constants. The orbit-counting and restriction arguments above
are the proofs in arbitrary dimension. A finite sample is not
used to extrapolate the result, and the published branching
estimate is not claimed as a new theorem of this run.
